\begin{abstract}
Modern maritime trajectory prediction is increasingly
\emph{context-aware}: conditioning on the surrounding environment
(coastline, channels) and social context (neighbouring vessels) reduces
displacement error by tens of metres. We observe that the cost of this
accuracy has migrated from the model to the \textbf{context-retrieval
pipeline}: for every AIS anchor, the pipeline issues an OpenStreetMap
(OSM) range query, computes a $128\!\times\!128$ signed-distance field
(SDF), and performs a $k$-nearest-neighbour lookup---millions of times
per corpus. Implemented by brute force in state-of-the-art systems, this
retrieval dominates the end-to-end training budget ($\approx\!17$
CPU-days per corpus), yet it has never been treated as a database
workload.

We present \textbf{M-CTX}, a spatial-indexing framework that recasts
context retrieval as a unified, indexable workload and accelerates it by
two to three orders of magnitude with provable correctness and
\textbf{bit-exact} downstream equivalence. At its core, \textbf{BR-LZ} is
the first one-curve learned spatial index to \emph{guarantee} recall
completeness for MBR-overlap range queries
(Theorem~\ref{thm:brlz-recall}), through a segment-local extent that
provably tightens candidate over-approximation by
$\mathbf{1.1\times}$--$\mathbf{2.7\times}$ over the global expansion
competitors require, at the fastest build and smallest footprint in its
class. M-CTX further contributes a linear-time SDF engine
($\mathbf{163\times}$ faster, numerically exact) and a storage--accuracy
co-design that compresses the context cache $\mathbf{64\times}$
($720$\,GB$\rightarrow$$11$\,GB) within $0.04$\,m ADE. Across four real
maritime regions, nine baseline systems, scaling to $40$\,M features, and
$10^{7}$-record streaming, M-CTX preserves the predictions of four
pretrained models byte-for-byte while cutting context construction from
$\approx\!11$ hours to under $3$ minutes.
\end{abstract}

\begin{IEEEkeywords}
spatial indexing, learned index, moving-object index, signed distance
field, AIS, trajectory analytics.
\end{IEEEkeywords}
